\section{Space Complexity --- Machines --- Sim}
%%% generated by the blueprint pipeline from TCSlib.Complexity.SpaceComplexity.Machines.Sim
\subsection*{Space Complexity --- Machines --- Sim: Introduction}

Declarations of the Lean module \texttt{TCSlib.Complexity.SpaceComplexity.Machines.Sim}.

\subsection*{Space Complexity --- Machines --- Sim: Results}

\begin{definition}[Track position]
\lean{Complexity.LogProg.trackPos}
\label{Complexity.LogProg.trackPos}
\leanok
\uses{Complexity.LogProg.Seg, Complexity.LogProg.TPos, Complexity.LogProg.wlen}
For a list of segments and a segment index $s$, define the track position of the left
boundary to be $-1$, the position of cell $c$ to be $c$, and the position of the right
boundary to be the length of segment $s$.
\end{definition}

\begin{definition}[Cell-position indicator]
\lean{Complexity.LogProg.TPos.isCell}
\label{Complexity.LogProg.TPos.isCell}
\leanok
\uses{Complexity.LogProg.TPos}
Define a Boolean indicator on track positions that is true for cell positions and false
for either boundary position.
\end{definition}

\begin{definition}[Consistency of a track position]
\lean{Complexity.LogProg.Consistent}
\label{Complexity.LogProg.Consistent}
\leanok
\uses{Complexity.LogProg.Seg, Complexity.LogProg.TPos, Complexity.LogProg.seg}
Let $S$ be a list of segments, let $s$ be a segment index, and let
$a,b\in\{\mathrm{false},\mathrm{true}\}$. Define a predicate on track positions by
requiring the left boundary to have direction $b=\mathrm{false}$ and the right boundary
to have direction $b=\mathrm{true}$. At a cell with parity $p$, require $p=a$ whenever
the flag of segment $s$ is true.
\end{definition}

\begin{lemma}[Agreement of abstract and tape-head moves]
\lean{Complexity.LogProg.gstep_tmove}
\label{Complexity.LogProg.gstep_tmove}
\leanok
\uses{Complexity.LogProg.Consistent, Complexity.LogProg.Seg, Complexity.LogProg.TPos, Complexity.LogProg.TPos.Valid, Complexity.LogProg.TPos.isCell, Complexity.LogProg.gstep, Complexity.LogProg.seg, Complexity.LogProg.tmove, Complexity.LogProg.trackPos, Complexity.LogProg.wlen}
\difficulty{6}
Let $S$ be a finite list of segments, let $s<	ext{length}(S)$, and let $t$ be a
tape-head position valid for segment $s$. Let $p,d	ext{ be Boolean flags such that }t$
is consistent with them: at the left boundary $d=	ext{false}$, at the right boundary
$d=	ext{true}$, and at a cell in a doubled segment its parity flag equals $p$. For any
movement $m\in\{-1,0,1\}$, apply the abstract update rule using segment $s$'s doubling
flag, whether $s$ is the last or first segment, and whether $t$ is a cell; denote its
new segment index, flags, and movement output by $s',p',d',m'$. Also apply the tape-head
move to obtain position $t'$. Then $s'$ is the segment index of $t'$, and $t'$ is
consistent with $p',d'$. Moreover, either $s'=s$ and the coordinate of $t'$ equals the
coordinate of $t$ plus $m'$, or $m'=0$ and the move crosses a segment boundary: from the
right boundary of segment $s$ to the left boundary of segment $s+1$, or from the left
boundary of segment $s$ to the right boundary of the preceding segment.
\end{lemma}

\begin{lemma}[Input symbol by head position]
\lean{Complexity.LogProg.inputSymbol_eq}
\label{Complexity.LogProg.inputSymbol_eq}
\leanok
\uses{Turing.Cfg, Turing.Cfg.inputSymbol, Turing.FinTM.inputSymbol_at}
\difficulty{4}
For any configuration on a Boolean input word $x$, the symbol under its input head
equals `none` when the head position is zero; otherwise it is the optional symbol of $x$
at index one less than the head position.
\end{lemma}

\begin{lemma}[Entries within the initial true run]
\lean{Complexity.LogProg.takeWhile_true_getElem?}
\leanok
\difficulty{4}
For every Boolean list $x$ and natural number $c$, if $c$ is less than the length of the
initial run of true entries in $x$, then the entry of $x$ at index $c$ is present and
true.
\end{lemma}

\begin{lemma}[Entry after the initial true run]
\lean{Complexity.LogProg.takeWhile_true_end}
\label{Complexity.LogProg.takeWhile_true_end}
\leanok
\difficulty{4}
For every Boolean list $x$, the entry immediately after its initial run of true entries
is either absent or false.
\end{lemma}

\begin{lemma}[Whole-Input Segment in Whole Mode]
\lean{Complexity.LogProg.Mode.seg0_whole}
\label{Complexity.LogProg.Mode.seg0_whole}
\leanok
\uses{Complexity.LogProg.Mode, Complexity.LogProg.Mode.seg0}
\difficulty{1}
Let $md$ be a mode equal to the whole-input mode, and let $x$ be a bit string. The
initial segment selected by $md$ from $x$ is $x$ itself, with its doubling flag set to
true.
\end{lemma}

\begin{lemma}[First-Component Segment in Unary Mode]
\lean{Complexity.LogProg.Mode.seg0_unaryFst}
\label{Complexity.LogProg.Mode.seg0_unaryFst}
\leanok
\uses{Complexity.LogProg.Mode, Complexity.LogProg.Mode.seg0}
\difficulty{1}
Let $md$ be a mode equal to the unary-first mode, and let $x$ be a bit string. The
initial segment selected by $md$ from $x$ is the leading run of $1$s in $x$, with its
doubling flag set to false.
\end{lemma}

\begin{lemma}[Initial segment of a call]
\lean{Complexity.LogProg.seg_callSegs_zero}
\label{Complexity.LogProg.seg_callSegs_zero}
\leanok
\uses{Complexity.LogProg.Mode.seg0, Complexity.LogProg.callSegs, Complexity.LogProg.seg}
\difficulty{1}
Let a call specification have a mode, let $x$ be a Boolean input, and let $W$ assign a
Boolean word to each register. Segment $0$ of the call's segment list is the initial
segment selected by the mode from $x$.
\end{lemma}

\begin{lemma}[Argument segments in a call]
\lean{Complexity.LogProg.seg_callSegs_succ}
\label{Complexity.LogProg.seg_callSegs_succ}
\leanok
\uses{Complexity.LogProg.callSegs, Complexity.LogProg.seg, Complexity.LogProg.seg_argSegs, Complexity.LogProg.seg_cons_succ}
\difficulty{3}
Let a call specification have an ordered argument-register list, let $x$ be a Boolean
input, and let $W$ assign a Boolean word to each register. For any index $a$ less than
the length of the argument list, segment $a+1$ in the call's segment list is the pair
consisting of the word assigned by $W$ to argument register $a$, and a flag that is true
exactly when $a+1$ is less than the length of the argument list.
\end{lemma}

\begin{lemma}[Doubling flag of a call segment]
\lean{Complexity.LogProg.segDbl_eq}
\label{Complexity.LogProg.segDbl_eq}
\leanok
\uses{Complexity.LogProg.callSegs, Complexity.LogProg.seg, Complexity.LogProg.segDbl, Complexity.LogProg.seg_callSegs_succ, Complexity.LogProg.seg_callSegs_zero}
\difficulty{4}
Let a call specification have a mode and an ordered list of argument registers, and let
$x$ be a Boolean input and $W$ assign a Boolean word to each register. For any index $s$
strictly less than one plus the number of argument registers, the flag indicating
whether segment $s$ is doubled equals the second component of segment $s$ in the list
formed from the mode's initial segment and the argument segments.
\end{lemma}

\begin{lemma}[Length of the first call segment]
\lean{Complexity.LogProg.wlen_zero_le}
\label{Complexity.LogProg.wlen_zero_le}
\leanok
\uses{Complexity.LogProg.Mode.seg0, Complexity.LogProg.callSegs, Complexity.LogProg.seg_callSegs_zero, Complexity.LogProg.wlen}
\difficulty{4}
For any call specification, input bit string $x$, and assignment $W$ of a bit string to
each register, the length of the first segment in the resulting call segments is at most
the length of $x$.
\end{lemma}

\begin{lemma}[Register of an argument segment]
\lean{Complexity.LogProg.segReg_eq}
\label{Complexity.LogProg.segReg_eq}
\leanok
\uses{Complexity.LogProg.segReg}
\difficulty{3}
Let a call specification have an ordered list of argument registers, and suppose there
is at least one register. For an index $s$ with $0<s\le$ the length of that list, the
register assigned to segment $s$ is the argument register at zero-based position $s-1$.
\end{lemma}

\begin{definition}[Work-tape register position]
\lean{Complexity.LogProg.regPos}
\label{Complexity.LogProg.regPos}
\leanok
\uses{Complexity.LogProg.TPos, Complexity.LogProg.callSegs, Complexity.LogProg.trackPos, Complexity.LogProg.wlen}
Let $c_0$ be a configuration with $m$ work tapes, let a call specification select an
ordered list of registers, let $W$ assign a bit string to each register, and let $s$ and
$t$ be a segment index and track position. For register $r$, its position is the
position in $c_0$ if $r$ is not selected. If it is selected, write $j$ for its index in
the selected list plus one: its position is the length of segment $j$ when $j<s$, is
$-1$ when $s<j$, and is the track position at segment $s$ when $s=j$.
\end{definition}

\begin{definition}[Input-head position]
\lean{Complexity.LogProg.inPos}
\label{Complexity.LogProg.inPos}
\leanok
\uses{Complexity.LogProg.CallSpec, Complexity.LogProg.TPos, Complexity.LogProg.callSegs, Complexity.LogProg.trackPos, Complexity.LogProg.wlen}
Let a call specification and register-word assignment determine the segments of a
virtual input, and let $s$ be a segment index with track position $t$. The input-head
position is the track position at segment $0$, shifted by one and converted to a natural
number, when $s=0$; when $s\ne0$, it is the length of segment $0$ plus one.
\end{definition}

\begin{definition}[Relation between simulated tracks]
\lean{Complexity.LogProg.TrackRel}
\label{Complexity.LogProg.TrackRel}
\leanok
\uses{Complexity.LogProg.CSt, Complexity.LogProg.Consistent, Complexity.LogProg.TPos, Complexity.LogProg.TPos.Valid, Complexity.LogProg.callSegs, Complexity.LogProg.inPos, Complexity.LogProg.regPos, Complexity.LogProg.vpos, Complexity.LogProg.vword, Turing.Cfg}
Let $c_0$ be a configuration with $m$ work tapes, and let $g$ and $d$ be configurations
of machines with respectively $m+k_D$ and $k_D$ work tapes, on the input and virtual
input determined by a call specification and register-word assignment. A track relation
at segment $s$, parity $a$, direction $b$, and track position $t$ requires that $s$ be
among the call's segments, that $t$ be valid and consistent there, and that the
simulated and virtual input-head positions agree with their prescribed positions. It
also requires the first $m$ work tapes and their head positions in $g$ to match $c_0$
and the prescribed register positions, the remaining tapes and head positions in $g$ to
match those of $d$, and the output of $g$ to equal the output of $c_0$.
\end{definition}

\begin{definition}[Simulation relation]
\lean{Complexity.LogProg.SimRel}
\label{Complexity.LogProg.SimRel}
\leanok
\uses{Complexity.LogProg.CSt, Complexity.LogProg.TPos, Complexity.LogProg.TrackRel, Complexity.LogProg.callSegs, Complexity.LogProg.vword, Turing.Cfg}
Let $g$ and $d$ be configurations as in the track relation, and fix a program node and a
call specification with its register-word assignment. They are related for simulation if
there exist a decider state $q$, a segment index $s$, Boolean parity and direction
values, and a track position $t$ such that $g$ is in the corresponding simulation state,
carrying the first bit of $d$'s output when present, $d$ is in state $q$, and the two
configurations satisfy the track relation at those parameters.
\end{definition}

\begin{definition}[Halting relation]
\lean{Complexity.LogProg.HaltRel}
\label{Complexity.LogProg.HaltRel}
\leanok
\uses{Complexity.LogProg.CSt, Complexity.LogProg.TPos, Complexity.LogProg.TrackRel, Complexity.LogProg.callSegs, Complexity.LogProg.vword, Turing.Cfg}
Let $g$ and $d$ be configurations as in the track relation, and fix a program node and a
call specification with its register-word assignment. They satisfy the halting relation
if there exist a segment index $s$, Boolean parity and direction values, and a track
position $t$ such that $g$ is in the return state carrying the first bit of $d$'s
output, or false if that output is empty; $d$ has halted; and the configurations satisfy
the track relation at those parameters.
\end{definition}

\begin{lemma}[Reading the Simulated Track]
\lean{Complexity.LogProg.track_read}
\label{Complexity.LogProg.track_read}
\leanok
\uses{Complexity.LogProg.CSt, Complexity.LogProg.Consistent, Complexity.LogProg.Mode.seg0, Complexity.LogProg.Mode.seg0_unaryFst, Complexity.LogProg.Mode.seg0_whole, Complexity.LogProg.TPos, Complexity.LogProg.TPos.Valid, Complexity.LogProg.TPos.isCell, Complexity.LogProg.TrackRel, Complexity.LogProg.callSegs, Complexity.LogProg.inPos, Complexity.LogProg.inputSymbol_eq, Complexity.LogProg.isCharRead, Complexity.LogProg.length_callSegs, Complexity.LogProg.regPos, Complexity.LogProg.segReg, Complexity.LogProg.segReg_eq, Complexity.LogProg.seg_callSegs_succ, Complexity.LogProg.seg_callSegs_zero, Complexity.LogProg.takeWhile_true_end, Complexity.LogProg.takeWhile_true_getElem?, Complexity.LogProg.trackPos, Complexity.LogProg.virtSym, Complexity.LogProg.vsym, Complexity.LogProg.vword, Complexity.LogProg.wlen, Complexity.LogProg.wlen_zero_le, Turing.Cfg, Turing.Cfg.inputSymbol, Turing.Cfg.workTapeSymbols, Turing.FinTM.bufferTape}
\difficulty{6}
Let $c_0$ be a configuration, let a call specification select a list of argument
registers, and let $W$ assign a Boolean word to each register. Assume each selected
register in $c_0$ stores its assigned word and the selected registers are distinct.
Suppose a configuration $g$ of the combined machine and a decider configuration $d$ are
related as follows: the current segment index is within the list of call segments, the
track position is valid and consistent with the parity and direction, the two input-head
positions represent the same virtual position, the original work tapes and their head
positions are preserved or positioned according to the track representation, the decider
tapes and head positions agree with $d$, and the output is unchanged. Define $r$ to be
the symbol read from $g$'s input head when the segment index is zero, and otherwise the
symbol under the corresponding argument-register head when such a register exists
(taking $r$ to be blank if there is no register). Then the machine reads a character
exactly when the track position is a cell, and the virtual symbol determined by the read
symbol, mode, segment, and direction equals the symbol represented by that track
position.
\end{lemma}

\begin{lemma}[Bounds on a Track Position]
\lean{Complexity.LogProg.trackPos_bounds}
\label{Complexity.LogProg.trackPos_bounds}
\leanok
\uses{Complexity.LogProg.Seg, Complexity.LogProg.TPos, Complexity.LogProg.TPos.Valid, Complexity.LogProg.trackPos, Complexity.LogProg.wlen}
\difficulty{3}
Let a list of segments and a segment index determine a segment length $L$. For any valid
track position, its associated integer position lies between $-1$ and $L$, inclusive.
\end{lemma}

\begin{lemma}[Register position after a track move]
\lean{Complexity.LogProg.regPos_after}
\label{Complexity.LogProg.regPos_after}
\leanok
\uses{Complexity.LogProg.TPos, Complexity.LogProg.callSegs, Complexity.LogProg.regPos, Complexity.LogProg.segReg, Complexity.LogProg.segReg_eq, Complexity.LogProg.trackPos}
\difficulty{6}
Let $c_0$ be a configuration with $m$ registers, let $x$ be a Boolean input, and let a
call specification with an ordered list of argument registers determine the virtual
input segments from $x$ and register words $W$. Assume the argument list has no
repetitions, $m>0$, and $s$ is at most the length of that list. For a register $r$, a
track position $tp$, a sign $mv\in\{-1,0,1\}$, and a second track position $T2$, suppose
either $G1=s$ and the track position of $T2$ equals that of $tp$ plus $mv$, or $mv=0$
and one of the following holds: $G1=s+1$, $tp$ is the right boundary and $T2$ the left
boundary; or $G1+1=s$, $tp$ is the left boundary and $T2$ the right boundary. Then the
register position at segment index $G1$ and track position $T2$ equals its position at
index $s$ and $tp$, plus $mv$ exactly when $s\ne0$ and $r$ is the argument register
selected at index $s$; otherwise the position is unchanged.
\end{lemma}

\begin{lemma}[Input head position after a segment move]
\lean{Complexity.LogProg.inPos_after}
\label{Complexity.LogProg.inPos_after}
\leanok
\uses{Complexity.LogProg.TPos, Complexity.LogProg.TPos.Valid, Complexity.LogProg.callSegs, Complexity.LogProg.inPos, Complexity.LogProg.moveInputPos_val, Complexity.LogProg.trackPos, Complexity.LogProg.trackPos_bounds, Complexity.LogProg.wlen_zero_le, Turing.moveInputPos, Turing.moveInputPos_zero}
\difficulty{5}
Let $x$ be an input word, let a call specification and register words determine a list
of segments, and let $s$ be a segment index. A track position is valid at an index when
it is a boundary or a cell within that segment. Suppose $tp$ is valid at $s$, $T_2$ is
valid at $G_1$, and either $G_1=s$ with the track coordinate of $T_2$ equal to that of
$tp$ plus the movement $mv$, or $mv=0$ and the positions cross between the right
boundary of segment $s$ and the left boundary of segment $s+1$, or between the left
boundary of segment $s$ and the right boundary of segment $s-1$. If the input head
position $p$ encodes track position $tp$ at segment $s$, then after moving the input
head by $mv$ when $s=0$ (and by zero otherwise), its position encodes $T_2$ at segment
$G_1$.
\end{lemma}

\begin{lemma}[Head of an appended optional bit]
\lean{Complexity.LogProg.head?_append_toList}
\leanok
\difficulty{3}
For every Boolean list $l$ and optional Boolean $o$, the head of the list obtained by
appending the elements of $o$ to $l$ is the head of $l$ if it exists, and otherwise is
$o$.
\end{lemma}

\begin{theorem}[One-step preservation of call simulation]
\lean{Complexity.LogProg.sim_step}
\label{Complexity.LogProg.sim_step}
\leanok
\uses{Complexity.LogProg.CSt, Complexity.LogProg.HaltRel, Complexity.LogProg.RProg, Complexity.LogProg.SimRel, Complexity.LogProg.TPos.Valid, Complexity.LogProg.TPos.isCell, Complexity.LogProg.TrackRel, Complexity.LogProg.callSegs, Complexity.LogProg.compileTM, Complexity.LogProg.ctr, Complexity.LogProg.gstep, Complexity.LogProg.gstep_tmove, Complexity.LogProg.head?_append_toList, Complexity.LogProg.inPos_after, Complexity.LogProg.inputSymbol_vpos, Complexity.LogProg.length_callSegs, Complexity.LogProg.regPos_after, Complexity.LogProg.seg, Complexity.LogProg.segDbl_eq, Complexity.LogProg.segReg, Complexity.LogProg.tmove, Complexity.LogProg.tmove_valid, Complexity.LogProg.toFin, Complexity.LogProg.track_read, Complexity.LogProg.vpos, Complexity.LogProg.vpos_tmove, Complexity.LogProg.vsym, Complexity.LogProg.vword, Turing.Action, Turing.Action.apply, Turing.Cfg, Turing.Cfg.inputSymbol, Turing.Cfg.workTapeSymbols, Turing.FinTM.bufferTape, Turing.MultiTapeTM, Turing.MultiTapeTM.step}
\difficulty{7}
Let $P$ be a program with $m$ registers and $d$ possible deciders, let $D$ be a decider
with $k_D$ work tapes and state set $S_D$, and choose an initial state for each decider.
Fix a call state $l$, a call specification $c_s$, an initial program configuration
$c_0$, and register words $W(r)$. Suppose $P$ makes the specified call at $l$, each
argument register's tape initially stores its corresponding word $W(r)$, and the
argument list has no repetitions. Suppose also that a combined configuration $g$ and
decider configuration $d_c$ are related as follows: the combined machine is in the
simulation state corresponding to $l$ and the decider's current state; the decider is
nonhalting; and for some valid position in the virtual input, the two input-head
positions, register tapes and heads, decider tapes and heads, and output agree with the
prescribed segment-by-segment simulation correspondence. After one step of $D$, if the
decider remains nonhalting, the resulting combined and decider configurations satisfy
the same simulation relation. If the decider halts, the resulting combined configuration
instead satisfies the return relation: the combined machine is ready to rewind the
input, the decider is halted, and the same tape, head-position, and output
correspondence holds.
\end{theorem}
